% TOP_PROBLEM: {"id": 3179, "title": "The Berge-Fulkerson conjecture", "category_id": 3, "category": {"id": 3, "name": "graph_theory", "display_name": "Graph Theory", "description": "Problems involving graphs, networks, and their properties.", "slug": "graph-theory", "order_index": 3, "created_at": "2026-07-31T15:26:25.671Z"}, "status": "open", "problem_number": "OPG-142", "set_id": 12}
% FIRST_POSTED: 2026-10-01
% ENTRY_KIND: statement_only
% UnsolvedMath ID 3179; source catalogue code OPG-142
% !TeX program = lualatex
% Definition and sources only; this file is not an LLM solution attempt.
\documentclass[11pt,a4paper]{article}
\usepackage[margin=25mm]{geometry}
\usepackage{fontspec}
\setmainfont{lmroman10-regular.otf}[BoldFont=lmroman10-bold.otf,
 ItalicFont=lmroman10-italic.otf,BoldItalicFont=lmroman10-bolditalic.otf]
\usepackage{amsmath,amssymb,mathtools}
\usepackage{unicode-math}
\setmathfont{latinmodern-math.otf}
\AtBeginDocument{\let\setminus\smallsetminus}
\usepackage{xurl,hyperref}
\hypersetup{colorlinks=true,urlcolor=blue}
\setcounter{secnumdepth}{0}
\setlength{\parindent}{0pt}
\setlength{\parskip}{5pt}
\setlength{\emergencystretch}{3em}
\begin{document}
\section*{The Berge-Fulkerson conjecture}\label{3179}
{\small UnsolvedMath ID 3179 · OPG-142 · Graph Theory · OpenGarden}
\subsection{Definitions and mathematical statement}
Let $G=(V,E)$ be a finite undirected loopless cubic
multigraph. Distinct parallel edges are retained, and
every vertex has degree three counting multiplicities.
Assume that $G$ is bridgeless, meaning that deleting any
edge does not increase its number of connected components.
A perfect matching is an edge set incident with every
vertex exactly once; in particular no two selected edges
have a common endpoint.

The Berge--Fulkerson conjecture asserts that for every
such $G$ there is a list of six perfect matchings
$M_1,\ldots,M_6$ satisfying
\[
 \sum_{i=1}^6\mathbf1_{\{e\in M_i\}}=2
 \qquad\text{for every }e\in E.
\]
The matchings in the list are not required to be distinct.
Their overlaps are permitted and indeed must achieve
exactly the displayed multiplicity on every edge.

Since $G$ is cubic, a perfect matching has $|V|/2$ edges
and $|E|=3|V|/2$. Thus the total edge appearances in six
perfect matchings equal $3|V|=2|E|$, as required, but
this counting identity alone does not ensure uniform
coverage. The conjecture requests that the matchings be
coordinated to give multiplicity exactly two, rather than
merely having union $E$.

When three perfect matchings partition $E$, repeating
each twice supplies the required list. This explains the
repetition convention and the role of the conjecture as
an extension of that special colouring situation.
\subsection{Short English statement}
Can six perfect matchings, allowing repeats, cover every edge of every bridgeless cubic graph exactly twice?
\subsection{Sources}
\begin{itemize}
\item[S1] ULAM.AI and the UnsolvedMath Contributors, supplied problems.json, record 3179 (OPG-142), OpenGarden collection. Catalogue identity and source transcription; CC BY 4.0.
\url{https://huggingface.co/datasets/ulamai/UnsolvedMath}.
\item[S2] Open Problem Garden, The Berge-Fulkerson conjecture. Original problem and discussion; the mathematical formulation below is expanded from the supplied source record.
\url{http://www.openproblemgarden.org/op/the_berge_fulkerson_conjecture}.
\end{itemize}
\end{document}
